\documentclass[11pt]{article}
\usepackage{amsmath, amssymb}
\usepackage{geometry}
\usepackage{graphicx}
\usepackage{hyperref}
\usepackage{listings}
\usepackage{xcolor}
\usepackage{booktabs}
\usepackage{caption}
\usepackage{subcaption}

\geometry{margin=1in}

\definecolor{codegray}{rgb}{0.95,0.95,0.95}
\lstset{
  backgroundcolor=\color{codegray},
  basicstyle=\ttfamily\small,
  breaklines=true,
  frame=single,
  language=Fortran,
  keywordstyle=\color{blue},
  commentstyle=\color{green!50!black},
}

\title{Correcting the Adjoint of Cubed-Sphere PPM Tracer Advection\\
in GCHP: Diagnosis and Fixes}
\author{Kay Suselj}
\date{July 2026}

\begin{document}
\maketitle

\begin{abstract}
Four-dimensional variational data assimilation (4D-Var) requires an
\emph{adjoint model} that propagates observation-error information backward
through time. A critical component is the adjoint of the tracer advection
scheme. In GCHP, tracer advection uses the Piecewise Parabolic Method (PPM)
on a cubed-sphere grid. We identified and corrected systematic errors in the
adjoint of the PPM boundary stencil at cubed-sphere face edges. These errors
caused the adjoint gradient to be inflated by a factor of 2--15 at face-edge
cells. This document describes the diagnosis, the two attempted fixes, and
the recommended corrected approach.
\end{abstract}

\tableofcontents
\newpage

%% ============================================================
\section{Background: Why the Adjoint Transport Matters}
%% ============================================================

In atmospheric CO$_2$ inverse modeling using 4D-Var, we seek surface flux
corrections $\delta\boldsymbol{\sigma}$ that minimize a cost function:
\begin{equation}
  J(\boldsymbol{\sigma}) = \frac{1}{2}
    \sum_{k}\left(\frac{y_k - H_k[\boldsymbol{\sigma}]}{\sigma_{\rm obs}}\right)^2
  + \frac{1}{2}\left(\frac{\boldsymbol{\sigma} - \boldsymbol{\sigma}_b}{\sigma_b}\right)^2,
  \label{eq:cost}
\end{equation}
where $y_k$ are satellite CO$_2$ observations, $H_k$ is the forward model
operator mapping fluxes to concentrations, and $\boldsymbol{\sigma}_b$ is a
background (prior) flux estimate.

The gradient of $J$ with respect to the control vector $\boldsymbol{\sigma}$
is computed by running the \emph{adjoint model}, which carries sensitivity
information $\lambda$ (units of CO$_2$ concentration per unit of flux)
backward through time. The adjoint of tracer transport---the dominant
operator in $H_k$---must therefore accurately propagate $\lambda$ from
observation locations back to source locations.

If the adjoint gradient $\nabla J$ is wrong by a factor of $\sim$2--15 at
face-edge cells (as we found), the gradient-descent optimizer (L-BFGS-B)
takes incorrect steps, causing flux fields to grow to unphysical values and
the cost function to oscillate rather than converge.

%% ============================================================
\section{The PPM Advection Scheme}
%% ============================================================

GCHP uses the Lin--Rood Finite-Volume dynamical core \citep{lin1996,putman2007},
which advects tracers with the Piecewise Parabolic Method (PPM). In the
simplified \emph{perfectly linear} variant (\texttt{mord=2}, \texttt{hord\_tr=2}),
the flux at the left edge of cell $i$ (with Courant number $c$) is:
\begin{equation}
  F_i =
  \begin{cases}
    q_{i-1} + (1-c)\bigl[a_i - q_{i-1} - c(a_{i-1}+a_i - 2q_{i-1})\bigr]
    & c > 0 \text{ (upwind from left)}, \\[6pt]
    q_i + (1+c)\bigl[a_i - q_i + c(a_i+a_{i+1} - 2q_i)\bigr]
    & c < 0 \text{ (upwind from right)},
  \end{cases}
  \label{eq:ppm_flux}
\end{equation}
where $q_i$ is the cell-mean tracer value in cell $i$ and $a_i$ is the
reconstructed tracer value at the \emph{left edge} of cell $i$, computed
from neighboring cell means:
\begin{equation}
  a_i = \frac{7}{12}(q_{i-1}+q_i) - \frac{1}{12}(q_{i-2}+q_{i+1})
  \quad \text{(interior cells)}.
  \label{eq:al_interior}
\end{equation}

%% ============================================================
\section{The Cubed-Sphere Grid and Ghost Cells}
%% ============================================================

\subsection{Grid structure}

The cubed-sphere grid partitions the sphere into six square faces, each
discretized on a $N_x \times N_y$ grid. In GCHP at resolution C24, each
face is $24 \times 24$ cells. The PPM stencil in Eq.~\eqref{eq:al_interior}
requires values at $i-2, i-1, i, i+1$. For cells near a face edge (e.g.\
$i=1$ or $i=N_x$), some of these neighbors lie on an \emph{adjacent face},
which is stored on a different MPI rank.

\subsection{Ghost cells (halo exchange)}

Before each PPM sweep, a \emph{halo exchange} (\texttt{mpp\_update\_domains})
copies data from neighboring MPI ranks into \emph{ghost cells}: the indices
$i = 0, -1, \ldots$ (left halo) and $i = N_x+1, N_x+2, \ldots$ (right halo).
These ghost cells hold valid, up-to-date field values and allow the PPM
stencil to be evaluated uniformly across face edges.

\subsection{Special boundary stencil for \texorpdfstring{$a_1$}{a(1)} and \texorpdfstring{$a_{N_x}$}{a(Nx)}}

At the face edges the standard 4-point stencil of Eq.~\eqref{eq:al_interior}
would extend two cells into the ghost region. Instead, the GCHP code uses a
special \emph{boundary stencil} that averages a ghost-side and an
interior-side second-order reconstruction:
\begin{equation}
  a_1 = \frac{1}{2}\underbrace{\left[
    \frac{(2\,\Delta x_0 + \Delta x_{-1})\,q_0 - \Delta x_0\,q_{-1}}
         {\Delta x_{-1} + \Delta x_0}
  \right]}_{\text{LEFT: ghost cells (indices 0, --1)}}
  + \frac{1}{2}\underbrace{\left[
    \frac{(2\,\Delta x_1 + \Delta x_2)\,q_1 - \Delta x_1\,q_2}
         {\Delta x_1 + \Delta x_2}
  \right]}_{\text{RIGHT: interior cells (indices 1, 2)}},
  \label{eq:al1_boundary}
\end{equation}
and symmetrically for $a_{N_x}$ at the right edge. The two remaining
stencil cells, $a_0$ and $a_{N_x+1}$, are computed from the ghost region
alone.

%% ============================================================
\section{The Reversed-Wind Adjoint Approach}
%% ============================================================

\subsection{Discrete transpose vs.\ reversed-wind adjoint}

The \emph{exact} adjoint of the discretized transport operator is its
matrix transpose. Computing this transpose requires either automatic
differentiation of the entire PPM code (expensive) or storing intermediate
states (checkpoint memory). An attractive approximation---used in GCHP---is
the \textbf{reversed-wind adjoint}: run the \emph{same} PPM code with
time reversed ($\Delta t \to -\Delta t$) and all wind fields negated. For
smooth, slowly-varying flows this approximates the true adjoint well, at
the cost of introducing errors proportional to flow nonlinearity and flux
limiter activity.

\subsection{Implication for Courant numbers at face boundaries}

Under reversed winds, a cell face that had $c > 0$ (inward flow, from ghost
to interior) in the forward pass now has $c < 0$ (outward flow, from
interior to ghost). This matters because:
\begin{itemize}
  \item In the \textbf{forward} with $c > 0$ at face $i=1$: upwind cell is
    the ghost cell $q_0$, and the flux depends on ghost-cell
    reconstructions $a_0$ and $a_1$ (half-ghost).
  \item In the \textbf{adjoint} with $c < 0$ at face $i=1$: upwind cell
    should be the interior cell $q_1$, and the flux should depend only on
    interior reconstructions.
\end{itemize}
The key problem is that neither the unmodified PPM code nor the first-attempt
fixes fully respect this distinction.

%% ============================================================
\section{Diagnosis: Finite-Difference Gradient Validation}
%% ============================================================

We validated the adjoint gradient against finite differences of the cost
function $J$. For each test cell $(f, j, i)$, we ran the forward model
twice with $\sigma \pm \varepsilon$ and computed
\begin{equation}
  g_{\rm FD} = \frac{J(\sigma+\varepsilon) - J(\sigma-\varepsilon)}{2\varepsilon},
  \label{eq:fd}
\end{equation}
and compared to the adjoint gradient $g_{\rm adj}$ at the same cell. For
a quadratic cost function and exact adjoint, $g_{\rm adj}/g_{\rm FD} = 1$.

\subsection{Baseline results (no adjoint fixes)}

Without any boundary stencil correction, cell ratios $g_{\rm adj}/g_{\rm FD}$
ranged from $-197$ to $+17$, with wrong signs at several face-edge cells.
This confirmed that the standard reversed-wind PPM adjoint is severely
incorrect at cubed-sphere face boundaries.

\subsection{Identifying the root cause}

Comparing cells at different positions from face edges revealed:
\begin{itemize}
  \item \textbf{Interior cells} (far from any face edge): ratio $\approx 1.1$,
    nearly correct.
  \item \textbf{Seam cells} (on one face edge, $i=0$ or $j=N_y-1$):
    ratio $\approx 2$--$5$.
  \item \textbf{Corner cells} (on two face edges simultaneously): ratio
    $\approx 15$.
\end{itemize}
The spatial pattern clearly pointed to the cubed-sphere face boundary
stencil as the source of error.

%% ============================================================
\section{Fix 1: MPI Adjoint Halo Exchange (Reverted)}
%% ============================================================

\subsection{Motivation}

After the PPM step, sensitivity values that accumulated in the ghost cells
(halo region) of the adjoint field $\lambda$ are never sent back to the
neighboring face that owns those cells. This ``lost'' sensitivity could
explain the large adjoint errors at face edges.

\subsection{Implementation}

We added a call to \texttt{mpp\_update\_domains\_ad} immediately after the
PPM advection step in \texttt{fv\_tracer2d.F90}:
\begin{lstlisting}
! After offline_tracer_advection in adjoint mode:
call mpp_update_domains_ad(q3, domain, ...)
\end{lstlisting}
This routine is the adjoint of the forward halo exchange: instead of
\emph{copying} ghost values from neighbors, it \emph{accumulates} ghost-cell
sensitivities into the owning cells on the neighboring face and zeros the
halo.

\subsection{Result and reason for reverting}

No improvement in gradient ratios was observed. After further analysis, we
determined that \texttt{mpp\_update\_domains\_ad} belongs to the
\textbf{discrete-transpose} adjoint framework, which is \emph{incompatible}
with the reversed-wind adjoint used in GCHP:
\begin{itemize}
  \item The reversed-wind adjoint propagates sensitivity by running the
    \emph{forward PPM} with negated winds. In this framework, the
    forward halo exchange already fills ghost cells with the correct
    sensitivity values (from the reversed-wind field on the neighboring
    face). Adding \texttt{mpp\_update\_domains\_ad} on top of this
    double-counts the cross-face contributions.
  \item The discrete-transpose approach requires storing all intermediate
    states and transposing every operator, which is not how the GCHP
    adjoint is implemented.
\end{itemize}
Fix 1 was reverted and is not re-applied.

%% ============================================================
\section{Fix 2: Ghost-Cell Boundary Stencil (Partial Improvement)}
%% ============================================================

\subsection{Motivation}

The forward PPM at $i=1$ with $c > 0$ (inward) uses ghost-cell
reconstructions $a_0$ and $a_1$ (half-ghost). Fix 2 hypothesized that the
adjoint at the same face with $c < 0$ (outward) should ``mirror'' this
by also using the ghost-cell stencil---symmetrically replacing the standard
interior stencil.

A runtime flag \texttt{tp\_core\_is\_adjoint} was added so that
\texttt{xppm}/\texttt{yppm} can switch to the boundary stencil only during
the adjoint run, leaving the forward model byte-identical.

\subsection{Implementation}

In the \texttt{mord=2} flux loop of \texttt{xppm}:
\begin{lstlisting}
#ifdef ADJOINT
  if (tp_core_is_adjoint .and. .not.nested .and. grid_type<3) then
    if (i==is .and. is==1 .and. xt < 0.) then
      qtmp = q1(i-1)          ! ghost cell q(0)
      flux(i,j) = qtmp + (1.+xt)*(al(i)-qtmp + xt*(al(i-1)+al(i)-(qtmp+qtmp)))
      ! uses al(0) [ghost] and al(1) [half-ghost]
    else if (i==ie+1 .and. (ie+1)==npx .and. xt > 0.) then
      qtmp = q1(i-1)          ! interior q(npx-1)
      flux(i,j) = qtmp + (1.-xt)*(al(i)-qtmp - xt*(al(i)+al(i+1)-(qtmp+qtmp)))
      ! uses al(npx) [half-ghost] and al(npx+1) [ghost]
    end if
  end if
#endif
\end{lstlisting}

\subsection{Results}

\begin{table}[h]
\centering
\caption{Finite-difference gradient ratio $g_{\rm adj}/g_{\rm FD}$ for Fix 2
  (5-day window, C24, hord\_tr=2, no convection/PBL).}
\begin{tabular}{llcc}
\toprule
Cell tag & Location & Position & Ratio \\
\midrule
\texttt{f0j03i00} & $25.6^\circ$S, $306.6^\circ$E & corner (x+y edge) & 14.66 \\
\texttt{f0j07i20} & $16.6^\circ$S, $16.1^\circ$E  & interior          & 0.71  \\
\texttt{f1j14i00} & $1.5^\circ$S,  $36.6^\circ$E  & x-seam ($i=0$)   & 2.06  \\
\texttt{f1j19i11} & $12.9^\circ$N, $22.3^\circ$E  & y-seam ($j=N_y$) & 2.06  \\
\texttt{f4j18i15} & $13.0^\circ$S, $286.1^\circ$E & interior (H-dominated)& 0.52 \\
Global uniform    & ---                            & ---               & 1.30  \\
\bottomrule
\end{tabular}
\label{tab:fix2_results}
\end{table}

The overall correlation improved dramatically from baseline (near zero) to
$r = 0.915$ (no conv/PBL) and $r = 0.866$ (with full physics). However, two
systematic errors remain:

\begin{enumerate}
  \item \textbf{Seam cells: ratio $\approx 2.06$.} The ghost-cell mirror
    stencil still propagates incorrect information at face edges. The root
    cause is identified below.
  \item \textbf{Corner cell: ratio $\approx 14.66$.} Compounding errors from
    both $x$- and $y$-direction face edges, exacerbated by the
    \texttt{copy\_corners} ghost operation.
\end{enumerate}

%% ============================================================
\section{Why Fix 2 Has a Factor-of-2 Error at Seam Cells}
%% ============================================================

Consider the left-face flux in adjoint mode ($i=1$, $c < 0$, outward flow):

\paragraph{Fix 2 formula:}
\begin{equation}
  F_1^{\rm Fix2} = q_0 + (1+c)\bigl[a_1 - q_0 + c(a_0 + a_1 - 2q_0)\bigr],
  \label{eq:fix2_left}
\end{equation}
where $q_0$, $a_0$ are ghost-cell values and $a_1$ is the half-ghost
boundary reconstruction from Eq.~\eqref{eq:al1_boundary}.

\paragraph{Correct interior-upwind formula:}
In the reversed-wind adjoint, $c < 0$ at face $i=1$ means flow is directed
\emph{outward} (from the interior toward the ghost region). The physical
upwind cell is therefore $q_1$ (interior, not ghost), and the reconstruction
should be computed from interior cells only. The correct formula is
\begin{equation}
  F_1^{\rm correct} = q_1 + (1+c)\bigl[a_1^{\rm loc} - q_1 +
    c(a_1^{\rm loc} + a_2 - 2q_1)\bigr],
  \label{eq:fix3_left}
\end{equation}
where
\begin{equation}
  a_1^{\rm loc} = \frac{(2\,\Delta x_1 + \Delta x_2)\,q_1 - \Delta x_1\,q_2}
                       {\Delta x_1 + \Delta x_2}
  \label{eq:al_loc_left}
\end{equation}
is the \textbf{interior-only half} of Eq.~\eqref{eq:al1_boundary} (the RIGHT
term alone, discarding the ghost LEFT term), and $a_2 = c_3 q_1 + c_2 q_2 +
c_1 q_3$ is already interior-only.

\paragraph{Why the ratio is $\approx 2$:}
Since $a_1 = \tfrac{1}{2}(A_{\rm ghost} + a_1^{\rm loc})$, Fix 2 effectively
uses a ghost+interior mixture where only the interior part is physically
meaningful for the adjoint. For smooth gradients where $A_{\rm ghost} \approx
a_1^{\rm loc}$ (which holds when the tracer field is symmetric across the
face), Fix 2's formula gives roughly $2 \times$ the correct interior
reconstruction---hence the observed ratio of $\approx 2.06$.

%% ============================================================
\section{Fix 3 (First Attempt, Incorrect): al\_loc as Half of al(1)}
%% ============================================================

\subsection{Idea}

Fix 3 attempted to correct the ratio-2 error by computing $a_1^{\rm loc}$
as the interior-only part of $a_1$. The implementation subtracted the ghost
half:
\begin{equation}
  a_1^{\rm loc,\,wrong} = 2\,a_1 - A_{\rm ghost}.
\end{equation}

\subsection{Error}

Unfortunately $a_1 = \tfrac{1}{2}(A_{\rm ghost} + a_1^{\rm loc})$ means
$2a_1 - A_{\rm ghost} = a_1^{\rm loc}$ is correct in principle, but the
implementation incorrectly computed $A_{\rm ghost}$ with wrong index offsets,
producing $a_1^{\rm loc,\,wrong} \approx 2\,a_1^{\rm loc}$---a factor-of-2
error in the other direction.

\subsection{Result}

The gradient correlation dropped from $r = 0.866$ (Fix 2) to $r = 0.711$
(Fix 3), with the slope falling from 0.97 to 0.46. Fix 3 was reverted.

%% ============================================================
\section{Fix 3 (Corrected): Interior-Only Boundary Reconstruction}
%% ============================================================

\subsection{Implementation}

Rather than subtracting the ghost contribution from $a_1$, compute
$a_1^{\rm loc}$ directly from interior cells using the same differencing
weights as the RIGHT term of Eq.~\eqref{eq:al1_boundary}:

\begin{equation}
  a_1^{\rm loc} = \frac{(2\,\Delta x_1 + \Delta x_2)\,q_1 - \Delta x_1\,q_2}
                       {\Delta x_1 + \Delta x_2},
  \label{eq:al_loc_correct_left}
\end{equation}
and for the right boundary:
\begin{equation}
  a_{N_x}^{\rm loc} = \frac{(2\,\Delta x_{N_x-1} + \Delta x_{N_x-2})\,q_{N_x-1}
                       - \Delta x_{N_x-1}\,q_{N_x-2}}
                      {\Delta x_{N_x-2} + \Delta x_{N_x-1}}.
  \label{eq:al_loc_correct_right}
\end{equation}

The complete adjoint flux formulas become:

\paragraph{Left face ($i=1$, $c < 0$, outward):}
\begin{equation}
  F_1 = q_1 + (1+c)\bigl[a_1^{\rm loc} - q_1 + c(a_1^{\rm loc} + a_2 - 2q_1)\bigr].
  \label{eq:fix3_correct_left}
\end{equation}

\paragraph{Right face ($i=N_x$, $c > 0$, outward):}
\begin{equation}
  F_{N_x} = q_{N_x-1} + (1-c)\bigl[a_{N_x}^{\rm loc} - q_{N_x-1}
    - c(a_{N_x-1} + a_{N_x}^{\rm loc} - 2q_{N_x-1})\bigr].
  \label{eq:fix3_correct_right}
\end{equation}

The same correction applies in the $y$-direction in \texttt{yppm}, replacing
$q_1, q_2, \Delta x_1, \Delta x_2$ with $q_{(i,1)}, q_{(i,2)},
\Delta y_1, \Delta y_2$ and so on.

\subsection{Key properties}

\begin{enumerate}
  \item \textbf{No ghost cells.} Equations \eqref{eq:fix3_correct_left} and
    \eqref{eq:fix3_correct_right} involve only indices $i \geq 1$ (interior)
    or $j \geq 1$ (interior). No MPI communication is required beyond what
    the standard forward halo exchange already provides for the reversed-wind
    adjoint sensitivity field.
  \item \textbf{Correct upwind cell.} For $c < 0$ (outward at left face), the
    upwind cell $q_1$ is interior. For $c > 0$ (outward at right face), the
    upwind cell $q_{N_x-1}$ is interior.
  \item \textbf{Consistent with reversed-wind framework.} The formula
    matches the standard mord=2 interior formula (Eq.~\ref{eq:ppm_flux},
    $c < 0$ branch) with $a_i$ replaced by $a_i^{\rm loc}$, which is the
    natural substitution when ghost cells are unavailable.
  \item \textbf{Forward model unchanged.} The fix is guarded by
    \texttt{\#ifdef ADJOINT} and the runtime flag \texttt{tp\_core\_is\_adjoint},
    so the forward PPM is byte-identical to the original.
\end{enumerate}

\subsection{Fortran implementation (xppm)}

\begin{lstlisting}
#ifdef ADJOINT
  if (tp_core_is_adjoint .and. .not.nested .and. grid_type<3) then
    if (i==is .and. is==1 .and. xt < 0.) then
      ! Left face outward: upwind = q1(1) = interior.
      ! al_loc = interior-only right half of al(1) formula.
      al_loc = ((2.*dxa(i,j)+dxa(i+1,j))*q1(i) - dxa(i,j)*q1(i+1)) &
             / (dxa(i,j) + dxa(i+1,j))
      qtmp = q1(i)
      flux(i,j) = qtmp + (1.+xt)*(al_loc - qtmp &
                + xt*(al_loc + al(i+1) - (qtmp+qtmp)))
    else if (i==ie+1 .and. (ie+1)==npx .and. xt > 0.) then
      ! Right face outward: upwind = q1(npx-1) = interior.
      ! al_loc = interior-only left half of al(npx) formula.
      al_loc = ((2.*dxa(i-1,j)+dxa(i-2,j))*q1(i-1) - dxa(i-1,j)*q1(i-2)) &
             / (dxa(i-2,j) + dxa(i-1,j))
      qtmp = q1(i-1)
      flux(i,j) = qtmp + (1.-xt)*(al_loc - qtmp &
                - xt*(al(i-1) + al_loc - (qtmp+qtmp)))
    else if (xt > 0.) then
      ...standard interior formula...
#endif
\end{lstlisting}

%% ============================================================
\section{Summary of All Changes}
%% ============================================================

\begin{table}[h]
\centering
\caption{Summary of adjoint advection fixes.}
\begin{tabular}{p{1.5cm} p{3.5cm} p{4.5cm} p{2.5cm}}
\toprule
Fix & Description & Outcome & Status \\
\midrule
Fix 1 & \texttt{mpp\_update\_domains\_ad} after PPM in \texttt{fv\_tracer2d.F90}
  & No improvement; double-counts cross-face sensitivity in reversed-wind framework
  & \textbf{Reverted} \\[6pt]
Fix 2 & Ghost-cell mirror stencil at face boundaries in \texttt{tp\_core.F90};
  added \texttt{tp\_core\_is\_adjoint} flag
  & Corr $= 0.915$ (no conv/PBL); seam ratio $\approx 2.06$; corner ratio $\approx 14.66$
  & \textbf{Current baseline} \\[6pt]
Fix 3a & Interior-only $a_1^{\rm loc}$ via wrong subtraction formula
  & Corr $= 0.711$ (worse than Fix 2); factor-of-2 error in $a_1^{\rm loc}$
  & \textbf{Reverted} \\[6pt]
Fix 3b & Interior-only $a_1^{\rm loc}$ from Eq.~\eqref{eq:al_loc_correct_left}--\eqref{eq:al_loc_correct_right};
  correct upwind cell
  & Pending validation; expected seam ratio $\to 1$
  & \textbf{Proposed} \\
\bottomrule
\end{tabular}
\label{tab:summary}
\end{table}

%% ============================================================
\section{Remaining Known Issues}
%% ============================================================

\subsection{Corner cells}

Cells at the intersection of two face edges ($i=0$ AND $j \approx 0$, for
example) show ratio $\approx 14.66$ even after Fix 2. These cells are
affected by three compounding sources:
\begin{enumerate}
  \item The $x$-direction seam stencil error (addressed by Fix 3b).
  \item The $y$-direction seam stencil error (also addressed by Fix 3b for \texttt{yppm}).
  \item The \texttt{copy\_corners} operation in GCHP, which fills corner ghost
    cells using rotational transforms. The adjoint of this operation is not
    currently implemented, leaving corner sensitivities unaccounted for.
\end{enumerate}
The corner cell fix requires additional work beyond the seam stencil and will
be addressed separately.

\subsection{Full-physics adjoint noise}

With convection and planetary boundary layer (PBL) mixing enabled, the
adjoint gradient becomes noisier because these processes are strongly
nonlinear and chaotic. Convective hot spots (e.g.\ the Maritime Continent)
amplify sensitivity values by factors of $\sim 10^3$ or more over a 5-day
adjoint run. This is a known limitation of the reversed-wind adjoint for
parameterized physics and is separate from the advection stencil issue.

%% ============================================================
\section{Practical Impact on 4D-Var Convergence}
%% ============================================================

The seam-cell gradient inflation ($\sim 2$) causes L-BFGS-B to take
$\sim 2\times$ too-large steps in the sigma (flux scaling) control variable
near face boundaries. Over several iterations this leads to:
\begin{itemize}
  \item Sigma values growing to unphysical magnitudes ($> 5$ in some OSSE runs).
  \item Cost function oscillation rather than monotonic decrease.
  \item Poor agreement between inferred and true fluxes in the OSSE validation.
\end{itemize}
After Fix 3b, the seam-cell ratios are expected to approach 1, which should
improve optimizer behavior significantly. The corner-cell error ($\sim 15$)
affects only a small number of cells but may still cause localized spikes in
the flux analysis.

%% ============================================================
\section*{Acknowledgments}
%% ============================================================

This work was carried out as part of the GCHP 4D-Var CO$_2$ OSSE at JPL,
with support from NASA.

%% ============================================================
\begin{thebibliography}{9}
\bibitem{lin1996}
  Lin, S.-J., and R.~B. Rood, 1996: Multidimensional flux-form
  semi-Lagrangian transport schemes. \textit{Mon. Wea. Rev.}, \textbf{124},
  2046--2070.

\bibitem{putman2007}
  Putman, W.~M., and S.-J. Lin, 2007: Finite-volume transport on various
  cubed-sphere grids. \textit{J. Comput. Phys.}, \textbf{227}, 55--78.
\end{thebibliography}

\end{document}
